\chapter{Il tranello della mensa}

\par\medskip
«Passiamo per la palestra e raggiungiamo le scale.»\\
«Va bene, Luca, fai tu strada.»

\par\medskip
I quattro imboccano la stessa scorciatoia che usavano le ragazze per scendere in mensa senza farsi vedere dai maschi quando il convitto era ancora attivo. La porta finestra di emergenza da cui sono entrati si apre infatti nell'ex camerata femminile e, da lì, un'uscita laterale conduce a un corridoio sul quale si affaccia una delle due porte della palestra. L'altra dà direttamente sul pianerottolo delle scale.

\par\medskip
Giovanni procede con cautela. Nella camerata ci sono ancora letti semidistrutti e materassi lasciati alla rinfusa sul pavimento, ostacoli da superare senza fare troppo rumore. Per fortuna Ipparchia e Rocky li seguono senza abbaiare, ma il loro obiettivo di segretezza sembra appeso a un filo.

\par\medskip
Luca gli afferra la mano.

\par\medskip
«Passiamo di qua.»

\par\medskip
Pochi passi ancora e raggiungono il corridoio.

\par\medskip
Una luce tremolante filtra dalla palestra. Proviene da un bidone di latta: uno stoppino fatto di stracci imbevuti di olio esausto continua a bruciare lentamente.

\par\medskip
Giovanni annusa. Conosce bene quell'odore. Ogni tanto scende dal secondo piano e gli pizzica il naso.

\par\medskip
Non gli piace.

\par\medskip
È sintetico, acre, tossico.

\par\medskip
Però fa luce.

\par\medskip
E la luce, a lui, non serve.

\par\medskip
\par\medskip
«Un attimo...»

\par\medskip
Luca controlla che tutti lo abbiano seguito, poi appoggia il cartone a terra per sistemarsi la cintura. Fatto questo, ripartono.

La porta della palestra è aperta. Anzi, la porta non c'è più da anni. Il varco però è inconfondibile: dove prima c'era un battente, ora si apre un portale di luce soffusa e suoni attenuati.

Dall'interno arrivano musichette elettroniche, jingle brevi e chiptune  di vecchie console. Anche quei suoni sono familiari a Giovanni, sorride alle frasi corte e simmetriche del Pac-Man, ripensando a quando studiava i percorsi di fuga seguiti dal ghost-team quando Pac-Man mangiava una delle pillole energetiche.

\par\medskip
Entrano con cautela. Nella palestra diversi ragazzini tengono in mano joystick e controller e giocano davanti a monitor recuperati.

\par\medskip
Luca li saluta con un cenno, ma non si ferma.

\par\medskip
L’attenzione di tutti è verso Rocky: la curiosità è alta, ma nessuno lascia la console, anche se gli occhi vanno e vengono dal cane.

\par\medskip
Attraversano l'ex palestra fino a raggiungere la seconda porta.  Nella mappa mentale che  Giovanni  sta tracciando, le due porte sono allineate creando una particolare connessione  che permette di entrare nella palestra da una porta e uscire dall'altra. Se  loro fossero in un videogioco di quelli in esecuzione nelle console, come il Pac-Man attraversandola si ritroverebbero al punto di partenza. Ma non è così, e quando attraversano il portale nel verso opposto  l'aura rosa azzurra che rifletteva suoi loro volti li abbandona e tornano ad essere  ombre scure nel convitto. Raggiungono il pianerottolo. Luca si avvicina all'amico: «Stai attento per le scale, ci sono diversi buchi».

\par\medskip
\begin{center}
  \includegraphics[width=1\linewidth]{piano 1.png}
\end{center}

\par\medskip
\begin{center}
  \includegraphics[width=1\linewidth]{piano 2.png}
\end{center}

\par\medskip
Giovanni si tiene radente al muro e misurando i passi uno ad uno raggiungono il primo piano.  Li la presenza delle casalinghe è palpabile o udibile due concetti analoghi quando l'udito è il primo dei sensi. il ritmo quasi isterico con cui gli oggetti vengono presi, puliti e riposti è una traccia inconfondibile. Tutti i piani sono soggetti alle loro cure, ma sempre uno alla volta a partire dal piano terra salendo fino al secondo. Comunque sono tutte nelle indaffarate nelle aulee, così  basta passare in un attimo: leggeri come il vento e continuare a scendere le scale.

\par\medskip
In mezzo alla rampa c’è un buco. Giovanni lo sa.  
È radente al muro perché conosce quel buco.

\par\medskip
Per un attimo pensa a Rocky, ma anche lui è e rimane radente al muro.

\par\medskip
Per ora nessun problema.\\
Un’altra rampa e raggiungono la Hall al piano terra.

\par\medskip
\begin{center}
  \includegraphics[width=1\linewidth]{piano terra.png}
\end{center}

\par\medskip
Qui Giovanni si muove sicuro, punta dritto alla porta della ex Mensa, mentre Luca lo segue con affanno.  
Il pacco comincia a farsi sentire.

\par\medskip
Rocky si guarda intorno. I cavalletti devono sembrargli ottimi per una pausa pipì, ma per fortuna non ne approfitta.

\par\medskip
Una candela è accesa. Qualcuno sta già mangiando.  
Ippa non cede al richiamo dell’odore, ma Rocky non è ancora altrettanto disciplinato.  
Si distacca dalla comitiva e va verso il tavolino.

\par\medskip
Luca lo guarda.\\
«No, Rocky, fermo!»

\par\medskip
È rapido. Non è ancora un ninja, ma ha la prontezza di lasciare andare i beni sostituibili per provare a salvare una vita non sostituibile.  
Il pacco cade mentre lui si lancia per afferrare il nuovo amico.

\par\medskip
Dopo sessanta secondi la situazione è questa: in mensa è rimasto solo Giovanni.  
Chi altro c’era si è dileguato.

\par\medskip
Luca è vigile e cosciente, ma malamente incastrato non riesce a far forza per risalire.  
Giovanni parla con lui mentre Rocky abbaia.

\par\medskip
«L'acqua mi arriva alla bocca. Non riesco a tenere su la testa.»\\
«Stai tranquillo, non salirà di livello.»\\
«Sì, sì.»\\
«Ci sono io, ti faccio uscire.»

\par\medskip
Giovanni chiama ancora aiuto, ma anche qui nel convitto le cose sembrano andare come fuori.

\par\medskip
«Vado a chiamare qualcuno. Torno subito.»\\
«Non lasciarmi da solo!»\\
«C’è Rocky, arrivo subito.»

\par\medskip
Ippa lo segue. Corre nella hall e poi verso l’entrata principale.  
Solleva la barra e si fa giusto lo spazio per passare.

\par\medskip
«Ehi amico, che passa?»\\
Una voce lo coglie appena si fa fuori.\\
«Ho bisogno. Uno dei ragazzini è finito nel pozzo.»
